\documentclass[11pt,reqno,a4paper]{amsart}
\usepackage[T1]{fontenc}
\usepackage{amssymb,booktabs}
\usepackage[hidelinks]{hyperref}
\allowdisplaybreaks[2]
\raggedbottom
\setlength{\emergencystretch}{1.8em}

\newtheorem{theorem}{Theorem}[section]
\newtheorem{proposition}[theorem]{Proposition}
\newtheorem{corollary}[theorem]{Corollary}
\newtheorem{lemma}[theorem]{Lemma}
\theoremstyle{remark}
\newtheorem{remark}[theorem]{Remark}
\newtheorem*{theoremT}{Theorem T}

\newcommand{\Fac}{\operatorname{Fac}}
\newcommand{\e}{\varepsilon}
\newcommand{\Y}{\mathcal Y}
\newcommand{\tern}[1]{\mathtt{#1}}
\newcommand{\R}{\mathbb R}
\newcommand{\Q}{\mathbb Q}
\newcommand{\Z}{\mathbb Z}
\newcommand{\dist}[1]{\lVert#1\rVert}
\newcommand{\rev}{\mathrm{rev}}
\newcommand{\word}[1]{\mathtt{#1}}
\newcommand{\HH}{\mathcal H}

\title[Squares and disjoint repeats in finite windows]
{Squares and disjoint repeats in finite windows with few palindromes}
\newsavebox{\affilbox}
\address{School of Artificial Intelligence, Shanghai Jiao Tong University, Shanghai, China}
\email{mqyqwerty@sjtu.edu.cn}
\date{}
\subjclass[2020]{11J68, 11J86, 68R15}
\keywords{Algebraic irrational numbers, integer-base expansions, squares, palindromes, combinatorics on words, critical exponent}
\hypersetup{pdfauthor={Qiyuan Ma},pdftitle={Squares and disjoint repeats in finite windows with few palindromes},pdfdisplaydoctitle=true}

\begin{document}
\sbox{\affilbox}{%
  \begin{minipage}{0.92\textwidth}
    \centering\normalfont\normalsize
    School of Artificial Intelligence,\\
    Shanghai Jiao Tong University, Shanghai, China\\
    ORCID: \href{https://orcid.org/0009-0000-9532-0821}{0009-0000-9532-0821}\\
    \href{mailto:mqyqwerty@sjtu.edu.cn}{mqyqwerty@sjtu.edu.cn}
  \end{minipage}}
\author[Qiyuan Ma]{Qiyuan Ma\\[6pt]\protect\usebox{\affilbox}}

\begin{abstract}
Bugeaud's Problem~10.42 asks whether every real algebraic irrational has, in every integer base, squares of unbounded root length and palindromes of unbounded length. We prove finite-window tests under palindrome restrictions. A ternary $9/4$-free word of length at least $64$ with at most six distinct palindromes contains a square whose root length is greater than $L/128$ and may be taken from an explicit sequence. Sufficiently long $16$-good or $17$-good ternary windows contain two disjoint equal blocks of length greater than $L/12000$. In base two, sufficiently long windows that are $13/5$-free with at most $18$ palindromes, or $28/11$-free with at most $20$ palindromes, contain a square of root length greater than $L/200000$. These combinatorial statements yield quantitative bounds on exceptional dyadic scales in a fixed algebraic expansion. We also record a ternary dichotomy for square starting positions, together with a local-frequency criterion that forces unbounded lengths when every maximizing component is a simple cycle. A branching construction shows that many square starts and a sublinear disjoint-repeat function need not force unbounded roots or palindromes. The two unbounded-length questions of Problem~10.42 remain open.
\end{abstract}
\maketitle

\section{Introduction}\label{sec:intro}

For an integer $b\ge2$ and a real number $x$, write $\{x\}=x-\lfloor x\rfloor$ and $\dist{x}=\min_{n\in\Z}|x-n|$. Let $\mathsf S_b(x)$ mean that the canonical base-$b$ expansion of the fractional part of $x$ contains factors $VV$ with $|V|$ unbounded. Let $\mathsf P_b(x)$ mean that it contains palindromes of unbounded length. Bugeaud's Problem~10.42 asks whether each property holds for every real algebraic irrational $x$ and every integer $b\ge2$ \cite[p.~220]{Bugeaud}. Adamczewski and Rampersad posed the corresponding questions in base two \cite{AR}.

The digit block starting after $i$ places and of length $d$ has integer value
\begin{equation}\label{eq:digitblock}
 D_{i,d}(x)=\lfloor b^{i+d}x\rfloor-b^d\lfloor b^i x\rfloor.
\end{equation}
Leading zeros are included. For a word $W$ let $[W]_b$ denote its integer value. Deleting any fixed prefix preserves both $\mathsf S_b$ and $\mathsf P_b$. A long palindrome can be trimmed symmetrically. If all unbounded squares began in a finite set of starts, one start would have unbounded roots. The second half of a sufficiently large square at that start contains each earlier given square at the same start, and moves that square beyond the deleted segment.

The long-repetition criterion of Adamczewski--Bugeaud--Luca, as stated by Bilu \cite[Theorem~3.3, p.~8]{Bilu}, implies that the greatest length of two disjoint equal factors among the first $N$ digits of a fixed real algebraic irrational is $o(N)$. The corresponding length $L_b(N)$ with distinct starts and possible overlap also satisfies
\begin{equation}\label{eq:known-early}
 L_b(N)=o(N).
\end{equation}
Indeed, overlapping matches of length $\ell$ at distance $d<\ell$ give a period-$d$ segment of length $\ell+d$. It contains $\lfloor\ell/d\rfloor+1$ adjacent copies of a length-$d$ word. Splitting them into two equal groups gives a disjoint match of length at least $\ell/3$, as in \cite[Lemma~3.4]{Bilu}.

Starting positions, distinct words and lengths are different quantities. Over a finite alphabet, unbounded square-root lengths are equivalent to infinitely many distinct squares. Infinitely many starts alone give less. Every infinite binary word has infinitely many square starts, because every length-four factor contains a square. The constructions of \cite{RSW,BC} can have a finite square catalogue. A recurrent palindrome of a prescribed length supplies information only at that length.

Adamczewski and Rampersad \cite[Theorem~3.1]{AR} proved that an algebraic ternary expansion has infinitely many square occurrences, or that $010$ or $02120$ is recurrent. The dichotomy recorded in Section~\ref{sec:ternary} concerns square starting positions and the recurrence of six specified words of length three. It is derived from established avoidance and transcendence results, using two explicit one-sided transfer lemmas; it gives no unbounded-length conclusion for square roots or palindromes.

The main results of the paper concern finite windows of the original digit string. The classification of infinite words with few distinct palindromes, due to Dvo\v r\'akov\'a, Mol and Ochem in the ternary case \cite{DMO} and to Dvo\v r\'akov\'a, Ochem and Opo\v censk\'a in the binary case \cite{DOO}, controls critical exponents of infinite words rather than proportional patterns in finite windows. Under explicit low-palindrome and power-free hypotheses, we show that a sufficiently long finite window contains a square or a pair of disjoint equal blocks whose length is a positive proportion of the window. Theorem~\ref{thm:six-pal} treats ternary $9/4$-free windows with at most six palindromes. Theorem~\ref{thm:ternary-linear-match} treats a broader class of $16$-good and $17$-good ternary windows, yielding a proportional disjoint repeat. Theorem~\ref{thm:binary-linear-square} gives the corresponding adjacent-square statement in base two. From these finite theorems we deduce bounds on the number of exceptional dyadic scales in a fixed algebraic expansion (Proposition~\ref{prop:dyadic-match}). The exceptional-scale count may be taken with exponent $3$ from the published quantitative Subspace Theorem of Evertse and Ferretti \cite{EF}, or with exponent $2$ from a $2026$ preprint interval form of that theorem \cite{Stephan}; both bounds are recorded.

Section~\ref{sec:frequency} gives exact criteria for an extremal local frequency to force unbounded lengths, together with a quantitative arithmetic deficit for algebraic paths, and branching constructions that delimit the criterion. In particular, Theorem~\ref{thm:branching-frequency} produces a positive-dimensional set of expansions for which the maximal disjoint equal-block length satisfies $R(N)=o(N)$, the quantity $D_n$ remains bounded, and both square-root lengths and palindrome lengths are bounded. Thus a large supply of square starts, even when combined with the known algebraic exclusion \eqref{eq:known-early}, does not by itself force the unbounded-length conclusions of Problem~10.42.

Both general questions---unbounded square-root lengths and unbounded palindrome lengths for every real algebraic irrational in every integer base---remain open.

\section{A ternary dichotomy from one-sided transfer}\label{sec:ternary}
The cuts prior to the five-letter iteration are finite prefix deletions.
\subsection{Statement and background}

For a finite or infinite word \(x\), a factor is a contiguous subword.
Write \(\Fac_m(x)\) for its set of length-\(m\) factors. A square is
\(UU\) with \(U\ne\e\), and a word is square-free if it has no square
factor. A factor of a one-sided infinite word is recurrent if its
set of starting positions is infinite. Positions in infinite words
are numbered from one.

\begin{theoremT}\label{thm:T}
Let \(\xi\) be a real algebraic irrational and let \(\mathbf a\) be the
canonical ternary digit word of its fractional part. Write
\[
P=\{\tern{010},\tern{020},\tern{101},\tern{121},\tern{202},\tern{212}\}.
\]
Then at least one of the following holds.
\begin{enumerate}
\item Infinitely many positions of \(\mathbf a\) are starting positions
of a square.
\item At least five distinct members of \(P\) are recurrent in
\(\mathbf a\).
\end{enumerate}
\end{theoremT}

Adamczewski and Rampersad \cite[Theorem 3.1]{AR} proved that an
algebraic ternary expansion has infinitely many square occurrences
or that \(\tern{010}\) or \(\tern{02120}\) is recurrent.
Theorem~T concerns square starting positions and recurrence of
specified words of length three. It gives no unbounded-length conclusion
for square roots or palindromes in Bugeaud's Problem~10.42 \cite{Bugeaud}.

Over a finite alphabet, infinitely many distinct squares are
equivalent to unbounded square-root lengths. Infinite binary words
with finitely many distinct squares exist \cite{BC}. They have
infinitely many square starting positions because every length-four
binary word contains a square.

Equivalently, the base-\(3\) value of a one-sided digit word is
transcendental if the word has a square-free suffix avoiding two
distinct members of \(P\). Indeed, finitely many square starting
positions and finitely many occurrences of the two words can all be
passed by one fixed cut. Such a suffix also excludes eventual
periodicity, so the value is irrational.

We derive Theorem~T as a corollary of established results. The two
one-sided lemmas in Section~\ref{sec:lifting} give the prefix control
needed to use the criterion of Adamczewski, Bugeaud, and Luca
\cite{ABL}.

\subsection{The avoidance inputs}

Put \(\Sigma_3=\{0,1,2\}\) and \(\Sigma_5=\{0,1,2,3,4\}\).
Badkobeh and Ochem \cite[Theorems 1 and 2]{BO} use the morphism
\(f=f_5\) defined by the images
\[
(f(0),f(1),f(2),f(3),f(4))
=(\tern{01},\tern{23},\tern{4},\tern{21},\tern{0}).
\]
Let \(t=f^\omega(0)\) and let
\[
\begin{split}
F_5=\{&
\tern{02},\tern{03},\tern{13},\tern{14},\tern{20},\tern{24},
\tern{31},\tern{32},\tern{40},\tern{41},\tern{43},\\
&\tern{121},\tern{212},\tern{304},\tern{3423},\tern{4234}\}.
\end{split}
\]
Let \(\Y\) consist of all one-sided infinite square-free words on
\(\Sigma_5\) that avoid \(F_5\). Theorem 1 gives \(t\in\Y\). Its proof
supplies the finite input
\begin{equation}\label{eq:window}
\begin{gathered}
pvs\text{ is square-free and avoids }F_5,\qquad
|p|=|s|=32,\quad |v|=8\\
\Longrightarrow\quad v\in\Fac_8(t).
\end{gathered}
\end{equation}
Section~\ref{sec:computation} describes an independent reproduction
of \eqref{eq:window}.

For \(i=1,2\) put
\[
B_1=\{\tern{010},\tern{020}\},\qquad
B_2=\{\tern{121},\tern{212}\}.
\]
Theorem 2 of \cite{BO}, with the definition of characterization
in that paper, asserts that every recurrent factor of any one-sided
infinite square-free word avoiding \(B_i\) occurs in \(s_i=M_i(t)\).
In the order \(0,1,2,3,4\), the images are
\[
M_1=(\tern{012},\tern{1},\tern{02},\tern{12},\e),\qquad
M_2=(\tern{02},\tern{1},\tern{0},\tern{12},\e).
\]
The recurrent languages also appear in Ochem and Rosenfeld
\cite[Theorems 2 and 3]{OR}, where the characterization is expressed
via distances between consecutive occurrences of a letter. Our use
of these classes requires statements about arbitrary one-sided words,
including their initial factors.

\subsection{Two one-sided transfer lemmas}\label{sec:lifting}

Badkobeh and Ochem \cite[\S3.1]{BO} give the four nonempty
images of \(M_i\circ f\), which appear as \(A,B,C,D\) below.
Define \(g_i=M_i\circ f^2\), so that \(s_i=g_i(t)\).
Direct composition gives the five images in the last column.
\[
\begin{array}{c@{\quad}cccc@{\quad}c}
\toprule
i & A & B & C & D & (g_i(0),g_i(1),g_i(2),g_i(3),g_i(4))\\
\midrule
1 & \tern{0121} & \tern{0212} & \tern{012} & \tern{021} & (AB,D,C,B,A)\\
2 & \tern{021} & \tern{012} & \tern{02} & \tern{01} & (AB,D,C,B,A)\\
\bottomrule
\end{array}
\]
Each of \(A,B,C,D\) has its only \(0\) at the beginning.
Every letter image under \(g_i\) has length between \(2\) and \(8\).

\begin{lemma}[Ternary suffix decoding]\label{lem:ternary}
If a one-sided infinite square-free word \(x\) on \(\Sigma_3\)
avoids \(B_i\), then \(x=H g_i(z_0)\) for a finite prefix \(H\)
and \(z_0\in\Y\). More precisely, if every length-\(41\) factor
starting at or after position \(N\) in \(x\) occurs in \(s_i\),
then \(H\) can be chosen with \(|H|\le N+10\).
\end{lemma}

\begin{proof}
Theorem 2 of \cite{BO} ensures that such an \(N\) exists. Each
length-\(41\) word outside \(\Fac_{41}(s_i)\) occurs only finitely
often in \(x\), and there are finitely many such words. Take \(N\)
one greater than the last starting position of any of them, or
take \(N=1\) if there is no such position.

Call the segment from one \(0\) to just before the next \(0\) a return
block. In a square-free ternary word every length-\(4\) factor contains
\(0\). A word of length \(4\) on \(\{1,2\}\) has an adjacent equal
pair or is an alternating square. Thus consecutive zero positions
in \(x\) differ by at most \(4\). If a return block \(R\) starts at
or after \(N\), the length-at-most-\(5\) word \(R0\) occurs in
\(s_i\) by right extension to length \(41\). All zero positions in
\(s_i=g_i(t)\) mark return-block starts by the table. The two zeros
of this occurrence of \(R0\) are consecutive, so \(R\) is one of
\(A,B,C,D\).

In the return-block sequence of \(g_i(t)\), an adjacent \(AB\)
occurrence is exactly the image of \(0\). For such a pair to cross
an image boundary, the image on the left would have to be \(g_i(4)\),
which ends in \(A\), and the one on the right would have to be
\(g_i(3)\), which starts in \(B\). The word \(t\) avoids \(\tern{43}\).
The local rule therefore pairs every adjacent \(AB\) and decodes it
as \(0\). It decodes the unpaired blocks \(D,C,B,A\) as \(1,2,3,4\),
respectively. All pairs are disjoint, and the neighbors of a block
determine whether it belongs to a pair.

Let \(q_0\) be the first zero position in \(x\) at or after \(N\), and
let \(q_1,q_2\) be the next two. Their bounds are \(q_0\le N+3\),
\(q_1\le N+7\), and \(q_2\le N+11\). Apply the pairing rule to the
return-block sequence beginning at \(q_0\). Cut at \(q_1\) unless
its block is the second member of a pair, in which case cut at
\(q_2\). The cut is an image boundary and has a preceding return
block starting at or after \(N\). This gives \(|H|\le N+10\)
and the exact factorization \(x=H g_i(z_0)\). The decoded word
\(z_0\) is infinite because each decoded letter uses at most two
return blocks.

To verify the legality of \(z_0\), consider any factor \(v\) of
length at most \(4\) in it. Its image \(g_i(v)\) uses at most eight
return blocks. Include the immediately preceding and following
return blocks in \(x\), and append the zero that marks the end
of the following block. The resulting segment has length at most
\(4(8+2)+1=41\). Its start lies at or after \(N\), including when
\(v\) is an initial factor of \(z_0\). By right extension to
length \(41\), the segment occurs in \(s_i\). The zeros identify
all of its return blocks there. The two neighbors guarantee the
same pairing decisions on the core. The occurrence in \(s_i\)
thus decodes to an occurrence of \(v\) in \(t\). All members of
\(F_5\) have length at most \(4\), so \(z_0\) avoids \(F_5\).
A square in \(z_0\) would map to a nonempty square in \(x\),
since \(g_i\) is nonerasing. Hence \(z_0\in\Y\).
\end{proof}

\begin{lemma}[Uniform five-letter desubstitution]\label{lem:five}
Every \(z\in\Y\) has the form \(z=p f(z')\), where \(p\) is a
prefix of length at most \(33\) and \(z'\in\Y\).
\end{lemma}

\begin{proof}
For every position \(s\ge33\) the length-\(8\) factor of \(z\)
starting at \(s\) has \(32\) letters on both sides in \(z\).
Since \(z\in\Y\), \eqref{eq:window} puts the factor in \(\Fac_8(t)\).

In the displayed images of \(f\), \(0,2,4\) occur exactly at image
starts, and \(1,3\) occur only as last letters of two-letter images.
Each \(1\) or \(3\) in \(t=f(t)\) is followed by an image start.
These are local rules on factors of length two, so position \(33\)
or \(34\) in \(z\) is an image start. From that position, read \(01\)
if \(0\) is followed by \(1\), and read the single letter \(0\)
otherwise. Read \(23\) or \(21\) at a \(2\), and read the single
letter \(4\) at a \(4\). These cases follow from the length-two
factors of \(t\) starting at those symbols. Each read ends just
before another image start. Thus the rule reads the whole suffix,
length \(32\) or \(33\) is discarded, and \(z=p f(z')\) for an
infinite word \(z'\).

The morphism \(f\) is nonerasing, so a square in \(z'\) would give
a square in \(z\). Suppose \(v\in F_5\) does occur in \(z'\).
Then \(f(v)\) begins at an image start at or after position \(33\)
in \(z\). For \(v\in F_5\) we have \(|f(v)|\le6\). This follows
from \(|v|\le3\) except at \(\tern{3423}\) and \(\tern{4234}\),
whose images have lengths \(6\) and \(5\), respectively.
Include the first symbol of the next image. The resulting
segment has length at most \(7\), and its right extension to
length \(8\) occurs in \(t\). Both its first symbol and the
appended symbol mark image starts in that occurrence in \(t\).
They fix the decoding, including at a terminal \(0\) of \(f(v)\).
Hence \(v\) occurs in the preimage \(t\) of \(t=f(t)\), a
contradiction. The right boundary matters because, for example,
\(f(\tern{14})=\tern{230}\) is a prefix of
\(f(\tern{10})=\tern{2301}\). In fact \(\tern{10}\) occurs
in \(f^4(0)\), so the bare image \(\tern{230}\) occurs in \(t\)
even though \(\tern{14}\) is forbidden. Therefore \(z'\) avoids
\(F_5\) and belongs to \(\Y\).
\end{proof}

\subsection{Prefix bounds and the arithmetic corollary}\label{sec:arithmetic}

Starting from \(z_0\) in Lemma~\ref{lem:ternary}, apply
Lemma~\ref{lem:five} successively. Its legality conclusion is the
hypothesis needed for the next application. Thus for every \(j\ge0\)
there are \(z_{j+1}\in\Y\) and a prefix \(p_j\) with
\[
z_j=p_j f(z_{j+1}),\qquad |p_j|\le33.
\]
In particular, for every \(k\ge1\) the original ternary word \(x\)
has the factorization
\begin{equation}\label{eq:iteration}
x=H\,g_i(p_0)\,g_i(f(p_1))\cdots
     g_i(f^{k-1}(p_{k-1}))\,g_i(f^k(z_k)).
\end{equation}

Write \(m_k=\min_{a\in\Sigma_5}|f^k(a)|\) and
\(L_k=\max_{a\in\Sigma_5}|f^k(a)|\). In the order \(0,1,2,3,4\),
the third iterates are
\[
(\tern{0123421},\tern{0423},\tern{01},\tern{0421},\tern{0123}).
\]
Every third iterate contains \(0\) and has length at least \(2\).
Since \(f^3(0)\) contains all five letters, every \(f^6(a)\) contains
all five letters. Also \(|f^6(a)|\le 64\). It follows that
\begin{equation}\label{eq:growth}
m_k\ge2^{\lfloor k/3\rfloor},\qquad L_{k+3}\ge2L_k,\qquad
L_k\le64L_{k-6}\le64m_k\quad(k\ge6).
\end{equation}
For the last inequality, use \(f^k(a)=f^{k-6}(f^6(a))\).
The sequence \(L_k\) is nondecreasing because \(f\) is nonerasing.
For \(1\le d\le k\) we have
\(L_{k-d}\le 2^{-\lfloor d/3\rfloor}L_k\), and thus
\begin{equation}\label{eq:sum}
\sum_{j=0}^{k-1} L_j
\le \left(2+3\sum_{h=1}^\infty 2^{-h}\right)L_k=5L_k.
\end{equation}
The part before \(g_i(f^k(z_k))\) in \eqref{eq:iteration} therefore
has length at most \(|H|+8\cdot33\cdot5L_k=|H|+1320L_k\).

Two of the first six letters of \(z_k\) are equal. Their images
under \(g_i\circ f^k\) are identical factors of \(x\). Let their
length be \(n_k\), their starting positions \(u_k<v_k\), and their
distance \(d_k=v_k-u_k\). Each letter image under \(g_i\circ f^k\)
has length between \(2m_k\) and \(8L_k\). There are at most five of
these images between the two starts. By \eqref{eq:growth}, for \(k\ge6\),
\begin{equation}\label{eq:dist}
\begin{aligned}
n_k &\ge2m_k,\qquad
d_k\le40L_k\le2560m_k\le1280n_k,\\
u_k-1 &\le |H|+1360L_k.
\end{aligned}
\end{equation}
The equality of the two factors supplies all comparisons needed
for period \(d_k\) on the length-\((d_k+n_k)\) segment starting at
\(u_k\). If \(d_k\le n_k\), this segment contains a square of root
length \(d_k\) at \(u_k\). Therefore \(d_k>n_k\).

Suppose now that \(x\) is a letter permutation of a suffix of a
ternary digit word \(\mathbf a=a_1a_2\cdots\), after discarding
\(\ell\) digits. The inverse permutation gives identical factors
at positions \(\ell+u_k,\ell+v_k\) in \(\mathbf a\). Let \(U_k\) be
the prefix ending at \(\ell+u_k-1\), and \(V_k\) the segment from
\(\ell+u_k\) to \(\ell+v_k-1\). Define \(V^w\) for \(w>1\) to be the
length-\(\lceil w|V|\rceil\) prefix of \(V^\omega\). Using
\eqref{eq:dist} with \(w=1+1/1280\), we have
\(\lceil d_k/1280\rceil\le n_k\).
The periodic segment in \(\mathbf a\) therefore gives
\begin{equation}\label{eq:abl}
\begin{aligned}
U_kV_k^w &\preceq \mathbf a,\qquad
          |V_k|=d_k>n_k\longrightarrow\infty,\\
\frac{|U_k|}{|V_k|}
&\le \frac{\ell+|H|+1360L_k}{d_k}
 \le \ell+|H|+43520.
\end{aligned}
\end{equation}
Here \(\preceq\) denotes the prefix relation and
\(L_k\le64m_k\le32n_k<32d_k\).
The transfer to \(\mathbf a\) uses only equality and length of
corresponding factors under the permutation.

\begin{proof}[Proof of Theorem T]
If both alternatives fail, choose two distinct members of \(P\)
with finitely many occurrences. Cut the digit word beyond their
last occurrences and beyond every square start. Encode a palindrome
\(aba\) by the arrow \(a\to b\). Under permutations of \(\Sigma_3\),
an unordered pair of distinct arrows has a common target, a common
source, opposite directions on the same edge, or forms a directed
path of length two. Representatives of the four types are
\[
\{\tern{010},\tern{212}\},\quad B_1,\quad B_2,\quad
\{\tern{010},\tern{121}\}.
\]

For \(\{\tern{010},\tern{121}\}\), the jump graph of Shallit and Shur
\cite[\S3.1, Figure 1 and Lemma 21(2)]{SS} excludes an infinite
square-free avoiding word. Every length-six factor of a square-free
word contains a jump \(aba\), since three distinct letters in every
consecutive triple would force a square of length six. Both
forbidden jumps belong to
\(\{\tern{010},\tern{121},\tern{202}\}\), one side of the graph.
With only \(\tern{202}\) left on that side, successive two-edge
returns give label pairs \(11,22\), or \(33\). The lemma excludes
an internal \(11\). A pair \(22\) followed by \(22\) gives the
forbidden \(222\) with a closed two-edge prefix. A pair \(22\)
followed by \(33\) gives an internal \(223\), also forbidden.
A tail of pairs \(33\) gives a forbidden \(333\) with a closed
two-edge prefix. An infinite walk is therefore impossible.

For \(\{\tern{010},\tern{212}\}\), the suffix also avoids
\(\tern{02120}\). The proof of \cite[Theorem 3.1]{AR} supplies
prefixes \(Y_k Z_k c_k Z_k\) on the normalized suffix, with
\(c_k\) a letter, \(|Y_k|\le2^{k+2}\) and \(|Z_k|\ge2^k/3\).
For all sufficiently large \(k\), the period word \(Z_kc_k\)
and the final \(Z_k\) give a \(3/2\)-power. After the inverse
permutation and the restoration of the fixed initial \(\ell\)
digits, the roots tend to infinity and the prefix-to-root ratios
are at most \(\ell+12\). For \(B_1\) and \(B_2\) the required
prefixes are those of \eqref{eq:abl}.

In all three cases with controlled prefixes, take a subsequence
on which the root lengths strictly increase. Since \(\xi\) is
irrational, its original digit word is not eventually periodic.
The criterion of Adamczewski, Bugeaud, and Luca
\cite[Theorem 1]{ABL}, applied in base \(3\) to that word, makes
\(\xi-\lfloor\xi\rfloor\) transcendental. This contradicts
algebraicity of \(\xi\).
\end{proof}

\subsection{Finite verification}\label{sec:computation}

The supplement's \texttt{verify\_window.py} enumerates the words
of length \(72\) on \(\Sigma_5\) that are square-free and avoid
\(F_5\). Starting with the empty word, it appends each letter.
When the prefix is legal, a new forbidden factor or square must
end at the appended letter. The checks are therefore all
forbidden suffixes and every square suffix with root length at
most half the new length. Induction on length establishes
completeness of this enumeration.

The enumeration gives \(1{,}012\) legal words of length \(72\)
and \(33\) distinct centers of length \(8\) at zero-based offsets
\(32\) through \(39\). Each center occurs in \(f^8(0)\), a prefix
of \(t\) of length \(72\). The JSON certificate records the
population at each length, one context witness and one prefix
offset per center, and a SHA-256 digest of the lexicographic stream
of legal length-\(72\) words. Running
\texttt{python verify\_window.py} beside the certificate repeats
the enumeration and checks the recorded data.

This computation verifies \eqref{eq:window}. Legality of the
arbitrary infinite preimages, the iteration, and the prefix
estimates come from the proofs of Lemmas~\ref{lem:ternary}
and~\ref{lem:five} and Section~\ref{sec:arithmetic}.


\subsection{Finite avoidance and a count of starts}\label{sec:finite}
Let $H(N)$ be the greatest length of a square-free factor of the first $N$ ternary digits that contains at most four of the six words in Theorem~T. Set $H(N)=0$ if the maximum is over the empty family. We record the finite-word input in a form that gives the loss estimates used below.

\begin{lemma}[Finite repeated factors]\label{lem:finite}
A square-free ternary word $W$ of length at least $1000$ that misses two of the six words has two equal length-$r$ factors at distance $d$, with
\[
 r\ge |W|/217600,\qquad r<d\le2560r .
\]
\end{lemma}

\begin{proof}
Associate the word $aba$ with the directed edge $a\to b$. The four types of pairs are a common endpoint, a common starting point, two opposite directions, and a directed path. Representatives are $\{010,212\}$, $\{010,020\}$, $\{121,212\}$, and $\{010,121\}$. The last class has an empty legal length-$24$ layer. For the first type use Hall desubstitution. For the middle two use the outer morphisms of the preceding section and then $f=f_5$. The finite-window verifier in the supplement gives the following local deconcatenation input, with legal preimages and both ends retained as losses. The outer cut has total loss at most $47$, with nonempty outer image lengths between $1$ and $8$. Above length $79$, each Hall or five-letter step strictly decreases the legal preimage length, keeps at least seven letters, and leaves two end words of total length at most $67$. The supplied proof verifies synchronization at both ends. The Hall fixed point is the one for $(012,02,1)$ \cite{Hall}.

For the inner morphism $\phi$, let $L_j$ and $m_j$ be the greatest and least lengths of the images of letters under $\phi^j$. The bounded windows are used only to establish the deconcatenation rules. The iteration is justified because each preimage remains legal. For $f$, the word $f^6(a)$ contains every letter and has length at most $64$, and all $f^3(a)$ have length at least $2$. For Hall the same statements hold with $3$ in place of $6$ and upper bound $27$ in the first assertion. Thus $L_j\le64m_j$, $L_{j+3}\ge2L_j$, and $\sum_{j<k}L_j\le5L_k$. Edge cases before the full incidence powers satisfy the same bounds by the image length estimates.

Stop when the inner word has length between $7$ and $79$. Iterating the upper bounds for the two end words gives
\[
 |W|\le47+8(67\cdot5+79)L_k\le3400L_k.
\]
Among its first six inner letters two coincide. Their original images have length at least $m_k$ and distance at most $40L_k$. A match at distance $d\le r$ would already contain a square with root length $d$. Permuting letters back preserves these lengths and positions.
\end{proof}

\begin{theorem}[Quantitative ternary windows]\label{thm:window}
For a real algebraic irrational $\alpha$, the function $H$ satisfies $H(N)=o(N)$. For $0<\eta\le1$ there is also the bound
\begin{equation}\label{eq:dyadic}
 \#\{j\ge1:H(2^j)\ge\eta2^j\}
 \ll_\alpha \eta^{-2}(1+\log(1/\eta))^2 .
\end{equation}
For any two of the six distinct words $P,Q$, the number of positions among the first $N$ that start a square, $P$, or $Q$ is
$\gg_\alpha(\log N)^{3/2}/\log\log N$ for sufficiently large $N$.
\end{theorem}

\begin{proof}
If $H(N)\ge\eta N$ for unbounded $N$ and fixed $\eta>0$, Lemma~\ref{lem:finite} produces two disjoint equal blocks of length at least $\eta N/217600$ in the first $N$ actual digits. The known long-repetition criterion \cite[Theorem~3.3]{Bilu} gives the contradiction. This deduction uses the finite desubstitution lemma to force the two blocks. The extra bound $d\le2560r$ in that lemma controls the exponent budget in the quantitative proof.

Here are the steps for \eqref{eq:dyadic}. Write $C=217600$, $K=2560$, and let $L\ge\eta N$ be the length of a bad factor, with $N=2^j$ and $L\ge4000$. Apply the lemma to its last $\lfloor L/3\rfloor$ letters and extend the two equal words leftward as far as their equality permits. They cannot extend to length $d$ inside the bad factor. Their preceding digits exist and differ, since their extension is shorter than $L/3$ and they started in the last third. If the first start is after $u$ digits, the second after $t=u+d$, and the new matching length is $\ell$, then
\[
 u\ge L/3,\quad d\ge L/(4C),\quad t\le N,\quad \ell\ge d/K .
\]
Let $p/(3^t-3^u)$ be obtained by extending that period forever after the first $u$ digits of $\alpha$'s fractional part. Then
\[
 |(3^t-3^u)\{\alpha\}-p|\le3^{-\ell},\qquad 3\nmid p .
\]
The divisibility assertion follows by comparing the unequal preceding digits of the two equal words. The height of $(3^t,3^u,p)$ is $3^t$.

Let $\theta=d/t$, split its possible range into dyadic bands $[s,2s]$, and split each band into $2K$ intervals. On one interval $[a,a']$, the normalized exponents for the real forms $X_1,X_2,\{\alpha\}(X_1-X_2)-X_3$ and the three coordinate forms at $3$ are bounded by
\[
 (0,-a,-1-a/K)\quad\hbox{and}\quad(-1,-1+a',0).
\]
The sum is $-3-\delta$, with $\delta=a/K-(a'-a)\ge s/(2K)$. The determinants have absolute value $1$. The interval form of the quantitative Subspace Theorem \cite[Theorem 3.3]{EF} gives one exceptional proper subspace and $O_\alpha(\delta^{-2}(1+\log(1/\delta)))$ height intervals. The large intervals have the form $[D,D^\omega)$ with $\omega\ll_\alpha\delta^{-1}$, and the remaining low height bound has $\log D_0\ll_\alpha1+\delta^{-1}$. Since $\eta2^j/3\le t\le2^j$, each large interval contains at most $O_\alpha(1+\log(1/\eta))$ of the selected dyadic indices. Summing over the bands, starting at $s\ge\eta/(8C)$, gives the right-hand side of \eqref{eq:dyadic}. Small $L$ and the low height range add at most logarithmically many indices.

For the exceptional subspace, split the dyadic indices modulo
\[
 g=\left\lceil\log_2(8C/\eta)\right\rceil .
\]
Along one class the selected $u,d$ both strictly increase. Three vectors are independent. Divide each row of their determinant by $3^u$. The two resulting product terms have distinct $3$-adic valuations $d_1-u_3$ and $d_1-u_2$, since $3\nmid p$. Each exceptional subspace thus contains at most two vectors per residue class, which contributes a smaller term.

Finally set $\eta=A\log J/\sqrt J$, for a sufficiently large fixed $A$. By \eqref{eq:dyadic}, a positive proportion of $j$ in $[J/2,J]$ satisfy $H(2^j)<\eta2^j$. In each good dyadic annulus take disjoint blocks of length $\lfloor\eta2^j\rfloor+1$. In each block there is a square or at least five of the designated three-letter words. In the latter case one of $P,Q$ occurs. The number of starts is $\gg J/\eta$. Taking $J=\lfloor\log_2 N\rfloor$ completes the count.
\end{proof}

The supplement gives the composite-base height normalization for factors using only three actual digits. It also gives the binary variants based on the Karhum\"aki--Shallit desubstitution and on the Badkobeh--Ochem short-square catalogues. These count occurrences, including finite root ranges in specified subclasses. The count does not improve to unbounded root lengths by a change of quantifiers. One elementary obstruction is to double each letter of a ternary square-free word. It has no $aba$ with $a\ne b$, and every square has root length $1$. An odd root of length at least $3$ would move a change of letter to the wrong parity, and an even root would give a square in the undoubled word.

\section{Transfer from finite windows with few palindromes}
\label{sec:few-palindromes}

A word is \(r\)-free if it has no factor of length at least \(rp\) and period \(p\) for an integer \(p\ge1\). The notation \(r^+\)-free allows equality. A palindrome catalogue includes the empty word. For an infinite word, \(D_r(N)\) counts the distinct starting positions of factors of exponent at least \(r\) that lie wholly in the first \(N\) digits. Write \(R(N)\) for the greatest length of two disjoint equal blocks in that prefix. A period need not be the least period.

\begin{lemma}\label{lem:pal-catalogue-cut}
A one-sided word with at most \(P\) recurrent nonempty palindromes has a suffix with at most \(P+1\) distinct palindromes, including the empty word.
\end{lemma}

\begin{proof}
If palindrome lengths were unbounded, one parity would occur with unbounded length. For every fixed smaller length of this parity one obtains palindromes at infinitely many distinct centers by symmetric trimming. The alphabet is finite, so at least one word of each such length recurs. This contradicts the bound \(P\). Once lengths are bounded, the finitely many nonrecurrent palindrome types can all be eliminated by taking a suffix past their last starts.
\end{proof}

\subsection{The ternary six-palindrome case}

\begin{theorem}\label{thm:six-pal}
A ternary \(9/4\)-free word of length \(L\ge64\), with at most six distinct palindromes, has a square of root length greater than \(L/128\). The root length can be chosen from
\[
                \ell_s=(4\cdot2^s-(-1)^s)/3,\qquad s\ge0.
\]
For each fixed \(s\), every sufficiently long word in this class has a square of root length \(\ell_s\).
\end{theorem}

\begin{proof}
The finite certificate in Round~64 gives \(780\) admissible length-\(19\) words. Each uses all three letters and has exactly two distinct letter doubles. The search appends one letter at a time, rejects a periodic suffix of length \(\lceil9p/4\rceil\) for any tested period \(p\), and accumulates the new palindromic suffixes. Independent enumeration tests all factors of the survivors and verifies the same word list. Consequently any \(W\) as in the theorem has three single-letter palindromes, two doubles and the empty word. These exhaust its catalogue. Name the letter without a double \(2\). A factor over \(\{0,1\}\) has length at most four, since a run of length three is forbidden and three runs would enclose a palindrome of length three or four at their middle run. An interval between consecutive \(2\)'s must consist of two such runs. It is \(0^a1^c\) or \(1^a0^c\) for \(a,c\in\{1,2\}\). Changing direction at a \(2\) gives a length-three palindrome. Swap the names \(0,1\) if necessary.

Discard at most five letters at each end, through the first and from the last \(2\). The remainder is \(E(v)\), where \(E\) inserts \(2\) into each pair \(10\) of a binary word and \(v\) is obtained by deleting the \(2\)'s. The pairs are disjoint, so \(n=|v|\ge2(L-10)/3\ge L/2\). Also \(v\) avoids \(000\) and \(111\). This comparison uses the insertion in the two-sided theorem of Dvo\v r\'akov\'a, Mol and Ochem \cite[Theorem~5]{DMO} with finite cuts.

If \(v\) has a factor of exponent at least \(7/3\) and period \(p\), use its prefix \(xyxyx\) with \(|x|=\lceil p/3\rceil\), \(|y|=p-|x|\). Its image is \(XYXYX\) with \(X=E(x)\) and \(Y\) including the internal insertions and at most two boundary insertions, with at most one when \(y\) is empty. If \(|x|\ge7\), cutting \(x\) at pairs \(10\) gives pieces of length at most four. Thus
\[
 |X|\ge5|x|/4-1,\qquad |Y|\le3|y|/2+2
                   \le3|x|+2<3|X|.
\]
The image has exponent greater than \(9/4\) at period \(|X|+|Y|\). For \(1\le|x|\le6\), the \(19{,}732\) choices of \(x,y\) satisfying \(|y|\le2|x|\) and giving \(xyxyx\) without constant triples are checked directly. In each image a factor has exponent at least \(9/4\). Of the nominal periods, \(53\) give equality and four pairs require finding a shorter-period factor. This proves that \(v\) is \(7/3\)-free.

Let \(\tau(0)=01,\tau(1)=10\). The finite-word decomposition of Karhum\"aki and Shallit \cite[Theorem~6]{KS} cuts at most two symbols at each end to obtain \(\tau\) of another \(7/3\)-free word. After \(k\) steps the length is at least \((n+4)/2^k-4>0\) for \(2^k\le n/4\). Take the largest such \(k\), which is at least three. This places some \(\tau^k(d)\) in \(v\). Each \(\tau^3(d)\) contains \(11\), giving two adjacent copies of \(V=\tau^{k-3}(1)\) in \(v\). The first digit of \(V\) is \(1\), so no insertion occurs at their seam and \(E(VV)=E(V)^2\). Its root has length at least \(2^{k-3}>n/64\ge L/128\). At every fixed smaller depth \(s\) with \(s+3\le k\), a \(\tau^{s+3}\)-image supplies the root \(E(\tau^s(1))\). To compute its length, the number \(C_s\) of changes in \(\tau^s(1)\) satisfies \(C_0=0\) and \(C_s=2^s-1-C_{s-1}\). The first digit is \(1\), and the last digit is \(1\) at even \(s\) and \(0\) at odd \(s\). It follows that the number of \(10\) pairs is \((2^s-(-1)^s)/3\), yielding \(\ell_s\).
\end{proof}

Both assumptions at the threshold matter. Infinite ternary \(7/4^+\)-free words exist by Dejean's theorem \cite{Dejean}. Dvo\v r\'akov\'a, Mol and Ochem \cite[Theorem~3(b),(c)]{DMO} give uniform codes of widths \(4,25\) from ternary \(7/4^+\)-free words. They yield respectively \(2^+\)-free words with at most seven palindromes and \(9/4^+\)-free words with at most six. The codes are injective and a complete image in a juxtaposition of two images starts only at a genuine cut, as verified at all offsets. A square of root length \(p\ge2q\), where \(q\) is the code width, therefore has \(p=mq\). The \(2m-1\) or more full images in it decode to a source block of period \(m\), of exponent at least \(2-1/m\). The source condition gives \(m\le4\). The respective root bounds are \(16,100\).

\subsection{Uniform exclusion of early disjoint matches}

\begin{proposition}\label{prop:dyadic-match}
For \(b=2,3\), a fixed real algebraic irrational \(\alpha\) and \(0<\epsilon<1/2\), the number of \(j\ge1\) for which \(R(2^j)\ge\epsilon 2^j\) is
\[
     O_{\alpha,b}\big(\epsilon^{-s}(1+\log\epsilon^{-1})^2\big).
\]
One may take \(s=3\) with the published theorem of Evertse and Ferretti, or \(s=2\) with the preprint interval theorem of Stephan.
\end{proposition}

\begin{proof}
Take \(\alpha\in(0,1)\) by subtracting an integer. If the equal blocks start after \(u,v\) digits with \(0\le u<v\) and have length \(h\ge\epsilon 2^j\), then \(v-u\ge h\) and the integer \(P=\lfloor b^v\alpha\rfloor-\lfloor b^u\alpha\rfloor\) satisfies
\[
 |e|=|(b^v-b^u)\alpha-P|<b^{-h}.
\]
Use the vector \((b^v,b^u,P)\), the forms \(X_1,X_2,\alpha(X_1-X_2)-X_3\) at the real place, and the coordinate forms at the prime \(b\). With \(K=b^v\), the respective exponent rows are \((1,u/v,-h/v)\) and \((-1,-u/v,0)\). Their total is at most \(-\epsilon\). The determinants have absolute value one. A real exponent is one, another is nonpositive, and the finite exponents are nonpositive with a zero. Thus Stephan \cite[Theorem~4.2]{Stephan} gives a bounded number of proper rational subspaces and \(O_\alpha(\epsilon^{-2}(1+\log\epsilon^{-1}))\) intervals for log heights, with upper-to-lower endpoint ratios \(O_\alpha(\epsilon^{-1})\), beyond the lower threshold of that theorem. Since \(\epsilon 2^j\log b\le\log K\le2^j\log b\), each interval covers \(O_\alpha(1+\log\epsilon^{-1})\) values of \(j\), as do the log heights below the threshold.

For a proper subspace take a plane \(aX_1+cX_2+dX_3=0\) containing it. If \(d=0\) then \(v-u\) is fixed, giving \(O(1+\log\epsilon^{-1})\) possible \(j\) by the same endpoint comparison. Otherwise the points \((b^{v-u},e/b^u)\) are collinear. Liouville's inequality gives a lower bound \(c_1b^{-C_1 2^j}\) on the absolute ordinate and the displayed error gives the upper bound \(b^{-\epsilon2^j}\). Partition the \(j\)'s into residue classes modulo \(\lceil C(1+\log\epsilon^{-1})\rceil\), with \(C\) fixed and sufficiently large. Successive points in each class have abscissa at least twice the previous one and absolute ordinate at most one quarter of the previous one, after excluding a bounded number of small \(j\)'s. Among three such points the second secant has at least the first horizontal increment and strictly smaller absolute vertical increment. Its slope cannot equal that of the first secant, even when the ordinates have different signs. Each plane therefore contributes \(O_\alpha(1+\log\epsilon^{-1})\) further scales.

For the published result, divide by the common power \(b^\sigma\) of the three coordinates to make the vector primitive. Since \(\sigma\le u\), the new exponents \(u',v'\) have \(v'-u'=v-u\) and the new log height still lies between \(\epsilon2^j\log b\) and \(2^j\log b\). The error decreases. Set \(k=\lceil4/\epsilon\rceil\) and split \(u'/v'\) into intervals \([l/k,(l+1)/k]\). For each class the normalized real exponents can be taken as \(0,(l+1)/k-1,-1-\epsilon\), and the finite exponents as \(-1,-l/k,0\). They are nonpositive with sum \(-3-\epsilon+1/k\). Evertse and Ferretti \cite[Theorem~3.3]{EF} give the corresponding fixed-weight interval bound and a proper subspace. The plane argument is unchanged since the ratios and the plane equation are preserved by division. The \(O(\epsilon^{-1})\) weight choices yield \(s=3\).
\end{proof}

\subsection{From window results to square lengths and counts of starts}

We record the counting argument once for the finite-window theorems used here. Suppose that every sufficiently long \(r\)-free window with the prescribed finite palindrome catalogue has two disjoint equal blocks of length \(>cL\), for a fixed \(c>0\). This holds if the window has a square of root length \(>cL\). Take the suffix supplied by Lemma~\ref{lem:pal-catalogue-cut}. In the first \(N\) digits, cut at each position counted by \(D_r(N)\), allowing a fixed initial cut at \(N_0\). The longest remaining consecutive window has length at least
\[
                    \frac{N-N_0-D_r(N)}{D_r(N)+1}.
\]
Each window is \(r\)-free. In a square-producing case, \(\liminf_ND_r(N)/N=0\) therefore forces actual unbounded square-root lengths, and each prescribed root supplied by arbitrarily long windows occurs arbitrarily late. If the square-root lengths are bounded, the theorem gives a fixed window length beyond which every late window has a full \(r\)-power. Disjoint windows then give \(D_r(N)\gg N\).

Assume now that the digits are those of an algebraic irrational in base two or three. In the separate shells \((2^{j-1},2^j]\), use only \(r\)-powers fully contained in the shell, and cut at their starts. Each remaining piece has length at most \(\max(L_0,c^{-1}R(2^j))\) for a fixed \(L_0\). If \(R(2^j)<\epsilon 2^j\), this shell gives \(\gg\epsilon^{-1}\) distinct starts for large \(j\) and sufficiently small \(\epsilon\). Put \(J=\lfloor\log_2 N\rfloor\). With \(\epsilon=A((\log J)^2/J)^{1/s}\) for sufficiently large fixed \(A\), Proposition~\ref{prop:dyadic-match} leaves a positive proportion of \(J/2\le j\le J\). Summing in distinct shells yields
\begin{equation}\label{eq:start-count}
        D_r(N)\gg_{\alpha}
        \frac{(\log N)^{1+1/s}}{(\log\log N)^{2/s}} .
\end{equation}
The implicit constant may also depend on the fixed combinatorial threshold. Using only \(R(N)=o(N)\) from the established long-repetition theorem \cite[Theorem~3.3]{Bilu}, every sufficiently late shell works for any fixed \(\epsilon>0\). Sending \(\epsilon\) to zero gives \(D_r(N)/\log N\to\infty\). Thus for at most five recurrent nonempty ternary palindromes, Theorem~\ref{thm:six-pal} yields \eqref{eq:start-count} with \(r=9/4\), as well as its unconditional symbolic implication from sparse starts. The next two sections give the other parameters.

\section{A four-letter descent for ternary windows}
\label{sec:ternary-descent}

We use the morphisms of Dvo\v r\'akov\'a, Mol and Ochem \cite[Section~6]{DMO}.
\[
\begin{array}{c|l|l}
a&\eta(a)&\gamma(a)\\\hline
0&010203&012021201210120102\\
1&2013&012101202120102\\
2&0132&012021201020121\\
3&0102013203&012021201210120212010201210120102
\end{array}
\]
Put \(u=\eta^\omega(0)\), \(v=\gamma(u)\), and let \(F_\eta\) consist of
\[
\begin{gathered}
 12,\ 21,\ 23,\ 31,\ 103,\ 302,\ 303,\ 132013,\ 320132,\\
 2010201,\ 2013201,\ 3013203,\ 030102030.
\end{gathered}
\]
A four-letter word is \(\eta\)-legal if it is square-free and avoids this list. A ternary word is called \(16\)-good if it is \(52/27\)-free and has at most \(16\) palindromes, and \(17\)-good if it is \(25/13\)-free and has at most \(17\). This section uses the stated strict thresholds.

\begin{lemma}\label{lem:first-gamma}
Every \(16\)-good or \(17\)-good word \(W\) of length \(L\ge4300\) contains, up to relabeling and reversal, a contiguous \(\gamma(z_0)\) with \(z_0\) \(\eta\)-legal and \(|z_0|\ge L/132\).
\end{lemma}

\begin{proof}
The first length for which no \(17\)-good word avoids \(010\) is \(132\), by complete enumeration. After relabeling this holds for every \(aba\) with \(a\ne b\). For \(L\ge134\), the interior of \(W\) has all six. Each is enclosed by the third letter \(c\) on both sides, as extending by \(a\) gives a double and extending by \(b\) gives an alternating square. Thus all six \(cabac\) occur. Together with the empty word and the single letters these account for \(16\) palindromes. Even palindromes are excluded by square-freeness. A palindrome of length at least nine would add two new lengths, seven and nine. In a length-seven palindrome the additional outer letter for the center \(cabac\) must be \(b\): \(c\) gives a double and \(a\) gives a relabeling of \(0102010\), which cannot be extended even by one letter while staying square-free. If the extra palindrome exists it can be named \(0120210\). Call the resulting prescribed catalogue the normalized \(17\)-good catalogue.

To give the finite checks a precise meaning, for each class and each window size \(d\) record all central length-\(d\) factors of admissible length-\(3d\) words. Treat these factors as edges with their length-\((d-1)\) prefix and suffix as vertices. After deleting \(d\) letters at both ends of an original word, its consecutive \(d\)-windows form a path in the graph. The reference language of \(v\) is computed exactly. Determine the source pairs of \(u\) by iteration until the rule taking internal \(\eta\)-pairs and seams of all current pairs stabilizes. Every letter occurs, so stability proves completeness. Once \(\min_a|\gamma\eta^k(a)|\ge d\), a reference \(d\)-window meets at most two adjacent images. Enumerating the factors of images of the allowed pairs gives precisely the reference set.

The checked graph for \(16\)-good words at \(d=36\) has \(318\) edges in two components of \(159\) each. Their labels are the \(36\)-factors of \(v\) and their reversals. The normalized \(17\)-good graph at \(d=36\) has \(884\) edges and \(518\) components, with \(566\) edges between components. The only internal edges give the same two label sets as in the \(16\)-good graph. A path can cross at most \(517\) times and traverse at most two stretches of internal edges. There are \(L-107\) windows after the end cuts. Hence there is a contiguous word \(U\) of length at least \((L-624)/2\), all of whose \(36\)-windows belong to one reference direction. There is no length-seven palindrome in those reference windows. Thus \(U\) is \(16\)-good, since \(25/13<52/27\). In the original \(16\)-good case take \(U=W\).

For \(16\)-good words at \(d=186\), the \(42{,}096\) admissible length-\(558\) words give \(1{,}776\) central words, exactly the \(888\) factors of \(v\) and their reversals. The corresponding vertex sets are disjoint. Delete \(186\) letters at each end of \(U\) and choose the reference direction by reversing if necessary. The marker \(M=\gamma(0)\) appears in each reference length-\(185\) word, and in \(v\) it occurs only at physical \(0\)-image starts. The latter assertion is checked on images of the \(12\) allowed source triples. The six first-return words to \(0\) in \(u\), with the next \(0\) excluded, are \(01,02,03,013,032,0132\), and their \(\gamma\)-images are distinct. These follow from the reference factors through length \(20\), since each reference \(19\)-window contains \(0\).

Select the first and last complete marker. Each consecutive marker pair together with the right marker fits in \(186\) positions and occurs in the reference. Indeed the \(185\)-window starting one digit after the left marker contains another complete marker whenever it fits. At the right boundary the claim is automatic from the remaining length. Thus the interval between the first and last marker starts parses into \(\gamma\)-images of source returns. Compared to \(U\), this trims at most \(372\) at each end. A square in the preimage would survive through the nonerasing map \(\gamma\), contradicting square-freeness of the window. Each \(\gamma(f)\) with \(f\in F_\eta\) has length at most \(186\) and has no occurrence, aligned or otherwise, in the reference language, by the same exact table. Hence the preimage is \(\eta\)-legal, and
\[
              |z_0|\ge (|U|-744)/33\ge(L-2112)/66\ge L/132 .
\]
\end{proof}

The stronger \(16\)-good input in the last step matters. Under just square-freeness and the displayed lists \(F_\eta,F_\gamma\) in \cite[Lemmas~6--7]{DMO}, the claimed central-word equalities fail. The exhaustive length-\(174\) and length-\(558\) enumerations respectively give \(5{,}548\) and \(1{,}446\) distinct centers, compared with \(274\) and \(888\) in the respective reference languages. For each discrepancy the supplement provides a concrete witness tested at all positions and a check of its nonmembership in the complete reference language. Lemma~\ref{lem:first-gamma} uses different hypotheses and establishes the first descent. A separate condition is needed for further descent.

\begin{theorem}\label{thm:ternary-linear-match}
A \(16\)-good or \(17\)-good ternary word of length \(L\ge40000\) has two disjoint equal blocks of length greater than \(L/12000\).
\end{theorem}

\begin{proof}
Write \(\Phi_k=\gamma\eta^k\). The factor \(\Phi_0(z_0)\) from Lemma~\ref{lem:first-gamma} is \(52/27\)-free. We maintain the pair of properties that \(z_k\) is \(\eta\)-legal and \(\Phi_k(z_k)\) is \(52/27\)-free.

Consider two symbol cuts \(s<t\) of \(z_k\), each with a next symbol, and let \(T_k=|\Phi_k(z_k[s:t])|\), with the right index excluded. Match leftwards from the cuts to the longest common suffix of their left contexts, and rightwards just through the common prefix of the images of the respective next letters. If these comparisons yield total matching length \(M_k\), then their union with the gap gives a factor of period \(T_k\) and length \(T_k+M_k\). Thus
\begin{equation}\label{eq:cut-ratio}27M_k<25T_k.\end{equation}
When the cuts are tested after appending a rightmost letter, all whole equal letters on the left contribute full image lengths. A first unequal pair contributes its image common suffix. No further matching is needed. For left extensions use the reverse procedure.

Extend the letter functions \(\lambda=(3,2,2,5)\), \(d=(-1,1,1,0)\) additively. Directly from the morphisms,
\[
\begin{gathered}
 \lambda(\eta(a))=6\lambda(a),\qquad d(\eta(a))=d(a),\\
 |\Phi_k(w)|=C_k\lambda(w)+(9/5)d(w),\qquad C_k=(33/5)6^k .
\end{gathered}
\]
The relative error from \(C_k\lambda(w)\) is at most \(\delta_k=3/(22\cdot6^k)\). Both tables of images are prefix and suffix codes in the sense that no image of one letter contains another image as a prefix or suffix. The property is preserved by composition. For distinct letters \(a,b\), denote the normalized limiting common prefix and suffix by \(P(a,b),S(a,b)\). If the corresponding common end of \(\eta(a),\eta(b)\) is \(J\), followed in the direction of the comparison by distinct \(a',b'\), the recurrence is \(6P(a,b)=\lambda(J)+P(a',b')\), with the analogous equation for \(S\). This gives
\[
\begin{array}{c|rrrrrr}
(a,b)&01&02&03&12&13&23\\\hline
 5P&1&6&11&1&1&6\\
 5S&5&0&10&0&5&0
\end{array}
\]
and the error after normalization by \(C_k\) is at most \(E_k=(5+2k)6^{-k}\). To see the bound, the actual comparison equals the full image of \(J\) plus the previous-level comparison at the unequal pair. After multiplication by \(6^k\), the error increases at most \(9|d(J)|/33\le2\) per step and initially is at most \(5\).

In the limiting comparison, let \(T\) sum the weights \(\lambda\) between cuts, and \(M\) sum those weights for full matching images and \(P,S\) for partial ones. If \(100M\ge93T\), then for \(k\ge6\)
\[
 \frac{M_k}{T_k}\ge
 \frac{(93/100)(1-\delta_k)-E_k}{1+\delta_k}>\frac{25}{27}.
\]
There are at most two partial comparisons and \(T\ge2\), which gives the error term \(E_k\) in the ratio. The last inequality holds at \(k=6\) by exact arithmetic and persists as the errors decrease. Hence later layers need only the necessary test \(100M<93T\). The first six layers use \eqref{eq:cut-ratio} exactly.

The finite descent certificate implements these seven versions of the test, on \(\eta\)-legal words. The longest legal word missing \(M_0=\eta(0)=010203\) has length \(31\). Consecutive complete markers together with the right marker thus fit in \(33\) symbols. The marker is unbordered. For each test the exhaustive enumeration of such intervals gives the same seven possible left-inclusive, right-exclusive return intervals. Six are the \(\eta\)-images of \(01,02,03,013,032,0132\), and the extra one is \(0102032\). Its interval followed by the right marker, and every \(\eta(f)\) for \(f\in F_\eta\) that passes the necessary test on its own, are the seeds for the next enumeration. Each test leaves six seeds of the latter kind. None of the seven seeds for any test extends by \(58\) symbols on \emph{both} sides while passing the test and the legality checks. This exhaustive assertion is implemented by extending right, then extending the reverse using the reversed forbidden list and interchanged prefix and suffix tables. The middle word, both surrounding parts and all inequalities are checked in the extended word.

If \(|z_k|\ge300\), choose full markers within the length-\(32\) bands starting \(58\) positions in from each end. Consecutive markers between the chosen two have the required two-sided padding, so only the six decodable intervals occur. This gives
\[
     z_k=P_k\eta(z_{k+1})Q_k,\qquad |P_k|,|Q_k|\le90 .
\]
The adjacent marker comparisons have genuine source returns. If \(z_{k+1}\) had a square, its nonerasing \(\eta\)-image would give a square in \(z_k\). If it contained \(f\in F_\eta\), the word \(\eta(f)\) would have \(58\) positions of padding on both sides in \(z_k\), which the seed test excludes. Hence \(z_{k+1}\) is legal. Its physical image \(\Phi_{k+1}(z_{k+1})=\Phi_k(\eta(z_{k+1}))\) is a factor of \(\Phi_k(z_k)\), giving the second maintained condition at the next layer.

Stop at \(12\le|z_k|\le299\). The lower bound follows from the last stage since the largest single \(\eta\)-image has length ten. Termination follows from minimum image length four. Write \(B_i=\Phi_i(z_i)\). As \(10\cdot6^i\le|\Phi_i(a)|\le35\cdot6^i\), the physical losses obey
\[
 |B_0|\le 299\cdot35\,6^k+180\cdot35\,\frac{6^k-1}{5}
 \le11725\,6^k .
\]
The first five source symbols include an equal pair. Their two complete images in \(B_k\) are disjoint and have length at least \(10\cdot6^k\). Since \(|B_0|\ge15L/132\), this length is greater than \(L/12000\). Undoing the comparison symmetries returns the two equal blocks to the original window.
\end{proof}

This descent also recovers the ternary language assertion in \cite[Theorem~8(b),(c)]{DMO}. On a two-sided \(16\)-good word, the first decoding and all subsequent levels have full padding. The allowed adjacent pairs at each level are in \(\mathcal L_2(u)=\{01,02,03,10,13,20,30,32\}\), as follows from the length-two forbidden words and square-freeness. At a large depth each fixed window crosses at most two source images, so it belongs to the language of \(v\). Conversely all reference factors occur, since every \(\eta\)-image contains \(0\), so each reference factor occurs in every sufficiently deep image. In a one-sided \(17\)-good word the \(36\)-graph path eventually stays in one reference component. Taking two-sided limits about occurrences arbitrarily far to the right identifies its recurrent language with the corresponding reference language. Here the original word need not be recurrent.

For a fixed algebraic irrational with at most \(15\) recurrent nonempty ternary palindromes use \(r=52/27\) in \eqref{eq:start-count}, and with at most \(16\) use \(r=25/13\). The result counts starts of powers with a threshold less than two. Under the latter palindrome bound there are also at least \(\lfloor(N-N_0)/450\rfloor\) starts of complete \(41/22\)-powers in the first \(N\) digits. This is a consequence of the published maximum \(449\) for \(41/22\)-free words with \(17\) palindromes \cite[Theorem~1(d)]{DMO}, also verified by the finite program. The two long equal blocks in Theorem~\ref{thm:ternary-linear-match} may be separated.

\subsection{A two-sided witness for the displayed four-letter condition}

For the Hall--Thue word \(T=\theta^\omega(0)\), with \(\theta=(012,02,1)\), use the nonerasing morphism
\[
\begin{array}{c|l}
0&0102030102013\\
1&010203201301020132030132\\
2&0102032013010201320301020130102030132 .
\end{array}
\]
Denote it by \(J\). The difference identity \(T_i=1+t_i-t_{i+1}\) for the binary Thue--Morse word \(t\) follows by comparing its four two-letter factors and placing \(\theta\)-cuts at \(2i+t_i\). In particular \(T\) avoids \(010\) and \(212\). A square in this difference word would give two matching binary difference blocks, so the corresponding two length-\((p+1)\) binary blocks differ everywhere by a constant, where \(p\) is the root length. Their shared endpoint forces the constant to be zero, giving an overlap in Thue--Morse, which is overlap-free \cite{KS}. This also proves square-freeness of \(T\).

The exact Hall--Thue language has \(24\) length-eight words. Every image of one of them under \(J\) is square-free, which suffices for square-freeness of \(J(T)\) by Currie's finite test \cite{CurrieTest}. Each checked image avoids \(F_\eta\), whose longest member has length nine, and the shortest \(J\)-image has length thirteen. Hence the checks also capture every possible occurrence of a forbidden factor. The word \(0102030102013010203\) appears in \(J(01)\) and is absent from the length-\(19\) reference language of \(u\). Primitivity of \(\theta\) ensures uniform recurrence of \(T\), and nonerasing \(J\) preserves it. Taking a two-sided limit about occurrences of the witness gives an \(\eta\)-legal two-sided word with a factor absent from \(u\). This contradicts the four-letter language assertion of the displayed Lemma~6 in \cite{DMO}.

The physical exponent constraint in the descent excludes this example. Thue--Morse contains squares \(VV\) with unbounded \(|V|\), by iterating \(\tau\) on \(11\). In the difference word they give \(ZaZ\), with \(|Z|=|V|-1\) and \(a\) the seam symbol. Under any fixed nonerasing morphism \(Q\), the exponents at the displayed period are at least \(1+|Q(Z)|/(|Q(Z)|+|Q(a)|)\to2\). Taking \(Q=\gamma\eta^k J\) works for every fixed \(k\), so this witness does not satisfy the inherited \(52/27\)-free condition.

\section{Adjacent squares in binary windows with few palindromes}
\label{sec:binary-finite}

Let the following morphisms be listed in source-letter order \(0,1,2\).
\[
  \varphi=(01,21,0),\qquad \mu=(011001,1001,0),\qquad \nu=(011,0,01).
\]
Write \(p=\varphi^\omega(0)\) and
\[
 F=\{00,11,22,20,212,0101,02102,121012,01021010,21021012102\}.
\]
Here an \(F\)-legal ternary word is cube-free and avoids \(F\). The morphisms, catalogue thresholds and the reference comparisons come from Dvo\v r\'akov\'a, Ochem and Opo\v censk\'a \cite{DOO}. The length estimates and the legal finite transfers below use the actual endpoints.

\begin{theorem}\label{thm:binary-linear-square}
If a binary word \(W\) has length \(L\ge150000\) and is either \(13/5\)-free with at most \(18\) palindromes or \(28/11\)-free with at most \(20\), it contains a square of root length greater than \(L/200000\). For the two respective cases \(g=\mu,\nu\), every sufficiently long such window has a square with root length \(|g\varphi^j(10)|\) for any fixed \(j\ge0\).
\end{theorem}

\begin{proof}
We give the finite inputs using the central-word digraph of Section~\ref{sec:ternary-descent}. The complete adjacent-pair set of \(p\) is \(01,02,10,12,21\), by closure under \(\varphi\). For a target length \(d\), the exact \(g(p)\)-language follows from these pair images at any depth with single-image length at least \(d\). The following table reports the checked central words and full windows.
\begin{center}
\begin{tabular}{@{}lrrr@{}}
\toprule
Window restriction & \(d\) & Full words & Central words\\
\midrule
\(13/5\)-free with at most \(18\) palindromes & 80 & 11324 & 644\\
\(28/11\)-free with at most \(20\) palindromes & 78 & 140320 & 2008\\
\(F\)-legal ternary words & 25 & 724 & 51\\
\bottomrule
\end{tabular}
\end{center}
In the first graph there are four components of \(161\) edges and no intercomponent edge. Their four label sets are \(\mu(p)\), its reversal, its binary complement and its complemented reversal, taken at length \(80\). In the second graph \(1336\) components have \(1384\) intercomponent edges. The only four components with internal edges form the \(156\)-edge reference sets of the four corresponding \(\nu(p)\) directions. The longest path in its acyclic component graph has length \(81\). The third central set is exactly the length-\(25\) language of \(p\).

In the first case cut \(80\) digits from each end. Apply the indicated comparison symmetry, so all remaining \(80\)-windows are factors of \(\mu(p)\). Every reference \(13\)-window contains the marker \(0101\), and a marker start plus two digits is a true \(\mu\)-cut. For each consecutive marker pair, the interval between the two cuts is one of
\[
   \word{0110011001}\mapsto01,\qquad
   \word{01001}\mapsto21,\qquad
   \word{011001}\mapsto0 .
\]
These are exact finite checks on the reference language. The distance from the first marker start to the next marker's end is at most \(14\), so checking the images of source length-\(14\) factors suffices. First and last marker cuts in the interior of \(W\) yield \(\mu(z_0)\) with a total loss at most \(190\). A source cube would lift to a cube in \(W\). If \(f\in F\) occurred in \(z_0\), take the factor spanning that occurrence and the bracketing complete markers. It has length at most \(38+20+4<80\), including the largest \(\mu(f)\), incomplete returns and markers. It occurs in \(\mu(p)\) and the same marker decoding would put \(f\) in \(p\), which avoids \(F\). Hence \(z_0\) is \(F\)-legal.

In the second case cut \(78\) at each end, giving \(L-233\) graph edges. At most \(81\) are intercomponent edges, and internal edges occur in at most four stretches. Thus a subword of length at least \((L-314)/4\) has all \(78\)-windows in a single reference direction. Work in the direction of \(\nu(p)\). Here every \(0\) is a cut, and the runs of \(1\)'s in \(W\) have length at most two, since \(W\) is cube-free. From first to last cut, the number of \(1\)'s between consecutive \(0\)'s decodes \(1,2,0\) for gaps with respectively zero, one and two such letters. Each end loses at most three digits. A cube lifts to \(W\). If \(z_0\) contains \(f\in F\), then \(\nu(f)0\) with the next marker has length below \(78\) and must occur in \(\nu(p)\), decoding \(f\) in \(p\). This proves legality also in the second case.

In both cases set \(B_0=g(z_0)\). We have \(|B_0|\ge L/5\) and \(|z_0|\ge L/30\). To descend, take any \(F\)-legal word \(z_i\) of length at least \(250\). The central \(25\)-window identity after cutting \(25\) at either end permits the same \(\varphi\)-decoding as in \(p\). Use the first and last \(0\) or \(2\) in that interval as cut positions. The pair \(21\) decodes as \(1\), and \(01\) decodes as \(0\). A \(0\) followed by \(0\) or \(2\) decodes as \(2\). The assumption of \(F\)-legality rules out the other possibilities. This gives
\[
                  z_i=P_i\varphi(z_{i+1})Q_i,\qquad
                  |P_i|,|Q_i|\le27.
\]
The surviving image is followed by a \(0\) or \(2\) guard. For \(f\in F\), its \(\varphi\)-image followed by the next cut letter has length at most \(19\) and would occur in the reference central language, at the same guarded cuts. Thus \(f\) cannot occur in \(z_{i+1}\). A cube in \(z_{i+1}\) would lift to a cube in \(z_i\). This preserves legality at every depth.

Stop at \(98\le|z_k|\le249\). The least length follows from the previous stage and the largest \(\varphi\)-image length two. The word shrinks at each stage, and \(|z_0|\le249\cdot2^k+54(2^k-1)<303\cdot2^k\), so \(k\ge5\). Put \(M_i=\max_a|g\varphi^i(a)|\), \(m_i=\min_a|g\varphi^i(a)|\). Every letter occurs in some \(\varphi\)-image, so \(M_i\) is nondecreasing. Checking the three letters gives \(|g\varphi^2(a)|\ge2|g(a)|\), whence \(M_{i+2}\ge2M_i\). Thus
\[
 \sum_{i<k}M_i\le3M_k,\qquad
 |B_0|\le249 M_k+54\sum_{i<k}M_i\le411 M_k .
\]
Every \(\varphi^3(a)\) contains all three letters and has length at most eight, giving \(M_i/m_i\le8\) for \(i\ge3\), with \(i=0,1,2\) checked directly. Every \(\varphi^5(a)\) contains \(1010\) and has length at most \(21\). Any letter of \(z_k\) now supplies an actual square
\[
                       \big(g\varphi^{k-5}(10)\big)^2
\]
in \(W\) under the comparison symmetry. Its root length is at least \(2m_{k-5}\ge M_k/84\), thus at least \(L/(5\cdot411\cdot84)>L/200000\). Increasing \(L\) forces \(k\) as large as desired, giving the fixed smaller scales too. The inverse of the comparison symmetry preserves literal equality of the two adjacent roots.
\end{proof}

The arXiv v2 proof of \cite[Theorem~7(a)]{DOO} uses a different length-\(20\) central graph for the first class. Exhaustive computation gives a single strongly connected component and \(282\) central words. For example,
\[
 \word{001001100101100100110010110011001001101001100110100110010011}
\]
is admissible and its central word \(\word{00101100110010011010}\) is outside the four complete \(20\)-letter reference sets. The length-\(80\) finite input avoids this discrepancy. Using the same decoding on both sides, all layers of an appropriate two-sided word stay in one reference direction, without end cuts. Every fixed factor appears within at most two adjacent full images at a large depth. All the recovered adjacent pairs belong to \(\mathcal L_2(p)\). Conversely, primitivity of \(\varphi\) puts each reference factor in every sufficiently deep image. This recovers the language assertion for the two-sided classes. The eventual single component on a one-sided path and two-sided limits identify its recurrent language as well.

For a binary digit string with at most \(17\) recurrent nonempty palindromes use \(r=13/5\) in Section~\ref{sec:few-palindromes}, and with at most \(19\) use \(r=28/11\). Sparse starts force unbounded square-root lengths in the original string. An algebraic irrational satisfies \eqref{eq:start-count} in the relevant case, and bounded roots would give a positive lower density of starts. The next construction shows that these two conclusions about starts are compatible with bounded roots.

\subsection{Equality at the exponent thresholds}

The \(21\)-uniform four-to-two-letter code \(H\) of Karhum\"aki and Shallit \cite[Lemma~8 and Fact~9]{KS} maps square-free words to \(7/3^+\)-free words. In source-letter order its images are
\[
 \begin{gathered}\word{011010011001001101001},\quad\word{100101100100110010110}, \\ \word{100101100110110010110},\quad\word{011010011011001101001}.\end{gathered}
\]
Let \(C_1,C_2\) be the \(72\)-uniform and \(49\)-uniform binary codes of \cite[Theorem~3(e),(f)]{DOO}. On \(7/3^+\)-free sources they give respectively \(13/5^+\)-free outputs with at most \(18\) palindromes and \(28/11^+\)-free outputs with at most \(19\). The images of \(C_1,C_2\) are provided by the verifier. Each is injective, and a complete \(C_i\)-image in the concatenation of two \(C_i\)-images starts only at a true boundary. A concatenation of two full \(H\)-images in \(H\) of any word also starts only at a true boundary. The two properties follow from the complete off-boundary checks with two and three neighboring images, respectively.

Take a four-letter \(7/5^+\)-free source \(a\), which is square-free. The composition \(y=C_i(H(a))\) has block width \(Q=21|C_i(0)|\). A match of length at least \(3Q\) aligns at full composite cuts. Indeed it contains two complete successive composite blocks from the first copy, \(C_i\)-synchronization aligns the internal cuts on the second copy, and the two decoded \(H\)-images fix the \(H\)-phase. A square of root length \(s\ge3Q\) thus has \(s=kQ\). At least \(2k-1\) whole composite images in it decode to a word of period \(k\) and exponent at least \(2-1/k>7/5\). Hence the roots are bounded by \(s<3Q\), giving \(4536\) and \(3087\) as strict bounds, respectively. All \(r\)-powers at the two thresholds have bounded periods since \(r>2\). Theorem~\ref{thm:binary-linear-square} gives a positive lower density of complete \(r\)-power starts, by partitioning into disjoint fixed-length windows.

\begin{lemma}\label{lem:entropy-repetitions}
In a finite-alphabet subshift of positive topological entropy \(h\), an ergodic measure of maximal entropy has \(R(N)=o(N)\) for almost every one-sided restriction. The subshift of four-letter \(7/5^+\)-free words has positive entropy.
\end{lemma}
\begin{proof}
In a finite-alphabet subshift the entropy map is upper semicontinuous, since measure entropy is the infimum of normalized finite digit entropies. The variational principle and the ergodic decomposition give an ergodic measure of entropy \(h\). For a fixed \(\epsilon>0\), a length-\(N\) word with a disjoint match of length \(s\ge\epsilon N\) is determined by the two starts, \(s\), and the prefix and suffix left after deleting the second copy. These last two pieces are admissible. The entropy estimate on the number of factors bounds the number of such length-\(N\) words by \(O_\delta(N^3)e^{(h+\delta)(1-\epsilon)N}\). The Shannon--McMillan--Breiman theorem \cite{Breiman} gives mass at most \(e^{-(h-\delta)N}\) for almost every point's prefix cylinder for all sufficiently large \(N\). Restrict to cylinders satisfying this mass bound and choose small \(\delta>0\). The probabilities are summable, so Borel--Cantelli gives the claim, using a countable sequence of \(\epsilon\) decreasing to zero.

Ochem \cite[Theorem~4.2(2)]{OchemSource} gives exponential growth for the finite four-letter \(7/5^+\)-free words. For each fixed length \(d\), a \(d\)-block with arbitrarily long admissible contexts on both sides is in the two-sided language, by compactness. The finitely many \(d\)-blocks therefore admit a uniform context size that guarantees this. Apart from two bounded ends, arbitrarily long finite words can be partitioned into consecutive such \(d\)-blocks. Their exponential growth rate is at most the number of two-sided \(d\)-blocks raised to \(1/d\). Letting \(d\) increase shows that the two-sided entropy equals the finite-word growth rate, and is positive.
\end{proof}

Using the almost-everywhere source from Lemma~\ref{lem:entropy-repetitions}, the composite synchronization lifts \(R_a(N)=o(N)\) to \(R_y(N)=o(N)\), since a disjoint match of length \(s\ge3Q\) gives a disjoint source match of length at least \(s/Q-2\). Restrict the source measure to a positive-measure set with the uniform cylinder bound \(e^{-(h-\delta)n}\) for all sufficiently large \(n\). A length-\(N\) output prefix fixes \(\lfloor N/Q\rfloor\) source digits. The image measure on binary cylinders has mass \(O(e^{-(h-\delta)N/Q})\). An interval of length \(2^{-N}\) meets at most three such cylinders. The mass distribution principle yields dimension at least \(h/(Q\log2)>0\) after taking \(\delta\downarrow0\). These are symbolic examples obeying the early-repetition constraint, with no assertion of algebraic irrationality.

\section{Local frequency bounds and the structure of maximizing cycles}
\label{sec:frequency}

We now work in any integer base \(b\). A finite list of forbidden factors is encoded by a digit graph \(G\) with memory \(m\ge1\). An edge with label \(a_0\cdots a_m\) goes from \(a_0\cdots a_{m-1}\) to \(a_1\cdots a_m\) and outputs \(a_m\). We retain only the allowed labels. A finite word on the path consists of the initial \(m\) digits and the outputs. Let \(w(e)\in\mathbb Q\) be an edge weight. Assuming \(G\) has a cycle, set
\[
  \rho=\max_C |C|^{-1}\sum_{e\in C} w(e),\qquad
  D_n=\rho n-\sum_{i=1}^{n}w(e_i).
\]
Here \(C\) ranges over the directed cycles and \(e_1,\ldots,e_n\) are the first \(n\) edges of the path. Write \(G_{\max}\) for the graph of edges on cycles of mean \(\rho\). A period word is called reversible up to rotation if its reversal is a cyclic rotation of it. For the estimates below a positive constant in a denominator is chosen large enough that \(D_n+C>0\) for every path.

\begin{theorem}\label{thm:extreme-frequency}
Suppose every strongly connected component of \(G_{\max}\) is a simple directed cycle. In the first \(m+n\) digits of any path, the maximum square-root length is at least
\[
                 c\frac{n}{D_n+C}-C
\]
for positive \(c,C\) depending only on \(G,w\). If every period word of these cycles is reversible up to rotation, the same assertion holds for the maximum palindrome length, with possibly different constants. In particular, \(\liminf_n D_n/n=0\) forces the respective unbounded lengths in the same digit string.
\end{theorem}

\begin{proof}
For each vertex \(v\), let \(h(v)\) be the largest sum of \(w(e)-\rho\) on a path ending at \(v\), with arbitrary start and with the empty path allowed. Removing closed subwalks cannot decrease this sum, so a path of length at most \(|V(G)|-1\) attains the maximum. For an edge \(e:u\to v\) put
\[
 c(e)=h(v)-h(u)+\rho-w(e)\ge0,
 \qquad
 \sum_{i=1}^n c(e_i)=D_n+h(v_n)-h(v_0).                         \tag{*}
 \label{eq:tight-cost}
\]
A cycle has mean \(\rho\) exactly when all its edges are tight, meaning \(c(e)=0\). Every internal edge in a cyclic strongly connected component of the tight graph lies on such a cycle, which identifies the cyclic components with those of \(G_{\max}\).

There are \(O_{G,w}(D_n+C_0)\) edges of positive cost on the path, by dividing the right side of \eqref{eq:tight-cost} by the least strictly positive cost if one exists. If no such edge exists there are zero to cut. Each resulting tight stretch crosses at most \(|V(G)|-1\) intercomponent edges, since the graph of components is acyclic. Cutting those leaves \(O_{G,w}(D_n+C_0)\) sections on the cycles and removes \(O_{G,w}(D_n+C_0)\) edges in total. One section has length at least \(c_1 n/(D_n+C_0)-C_1\).

A length-\(s\) section of a cycle of period \(p\) contains a square of root length \(p\lfloor s/(2p)\rfloor\) when \(s\ge2p\). The cycle periods are bounded. For a reversible period there is an integer \(a\) for which the cyclic digits satisfy \(d_i=d_{a-i\bmod p}\). Choosing the nearest center from those with \(2c\equiv a\bmod p\) and trimming symmetrically gives a palindrome of length at least \(s-2p\) when this lower bound is positive. This proves both estimates.
\end{proof}

\begin{theorem}\label{thm:extreme-arithmetic}
Under the simple-cycle hypothesis of Theorem~\ref{thm:extreme-frequency}, let a canonical expansion of an algebraic irrational \(\alpha\) follow the graph, possibly after a fixed prefix. Then, for all sufficiently large \(n\),
\[
                 D_n\gg_{\alpha,G,w}
                 \frac{(\log n)^{3/2}}{\sqrt{\log\log n}} .
\]
Here the count starts at the chosen suffix, and \(b\) is encoded by the graph's alphabet.
\end{theorem}

\begin{proof}
Let \(P\) be the least common multiple of the lengths of the simple cycles in \(G_{\max}\). Write \(x=0.a_1a_2\cdots\) in base \(b\) for the chosen suffix. As \(x\) is obtained by multiplying \(\alpha\) by a power of \(b\) and subtracting an integer, it is algebraic irrational. With \(N=m+n\) set
\[
 E_P(N-P)=\#\{1\le i\le N-P:a_i\ne a_{i+P}\}.
\]
A pair in the same periodic section of the preceding proof agrees. The initial \(m\) symbols, cut edges and last \(P\) positions of each section account for all possible mismatches. Thus \(E_P(N-P)\ll_{G,w}D_n+C_0\).

The fractional part \(y=\{(b^P-1)x\}\) is also algebraic irrational. It is the fractional part of \(\sum_{i\ge1}(a_{i+P}-a_i)b^{-i}\). For a length-\(s\) zero run of these signed differences, multiply by \(b^{i-1}\) for its initial index \(i\) and discard the integral contribution from earlier differences. The remaining signed tail is nonzero and has absolute value strictly less than \(b^{-s}\). Irrationality excludes equality and zero. Its base-\(b\) fractional digits start with \(s\) zeros or \(s\) copies of \(b-1\), according to the sign. Partitioning the first \(N-P\) positions into these zero runs and singleton nonzero differences shows that the number of adjacent digit changes in the corresponding prefix of \(y\) is at most \(2E_P(N-P)+O(1)\). Bugeaud and Evertse's digit-change bound \cite[Theorem~3.1]{BE} gives the asserted inequality. A fixed cut changes the prefactor \(b^P-1\) on \(\alpha\) by a fixed power of \(b\).
\end{proof}

\subsection{Branching and bounded lengths at the maximum}

We use the following binary source. Let \(a\) be a four-letter \(7/5^+\)-free word chosen from the positive-entropy subshift in Lemma~\ref{lem:entropy-repetitions}, and put \(z=H(a)\), with the 21-uniform code \(H\) from Section~\ref{sec:binary-finite}. A match of length at least \(63\) in \(z\) contains two consecutive full \(H\)-images and therefore aligns modulo \(21\) by the finite synchronizing check there. In particular, \(R_z(N)=o(N)\) when \(R_a(N)=o(N)\). Also, \(z\) has no factor of exponent at least \(9/5\) with period \(p\ge126\). A periodic factor of that length would have a match of length at least \((4/5)p\ge63\), so \(p=21d\) with \(d\ge6\). The full images inside the factor would give a source factor of period \(d\) and exponent at least \(9/5-2/d>7/5\). The same source measure from Lemma~\ref{lem:entropy-repetitions} is used for the dimension estimates.

\begin{theorem}\label{thm:branching-frequency}
If a strongly connected component of \(G_{\max}\) has a vertex with two outgoing edges in the component, there is a positive-Hausdorff-dimension set of real numbers with digit strings following \(G\) for which \(D_n=O(1)\), square-root lengths and palindrome lengths are both bounded, and \(R(N)=o(N)\).
\end{theorem}

\begin{proof}
Fix a branching vertex \(v\) in the component and output words \(U_0,U_1\) of two first returns from its distinct outgoing edges. The words have distinct first digits. The closed paths with words \(U_0U_1\) and \(U_1U_0\) have a common output length \(q\). Their \(n\)-fold concatenations are \(2^n\) distinct cyclic words with designated initial position, all of length \(nq\). Given one reflection \(i\mapsto a-i\bmod nq\), any word fixed by it is determined by a contiguous cyclic arc of length at most \(\lceil nq/2\rceil+1\). Such an arc contains at least one representative of every reflection orbit, by taking positions between the reflection axes and including adjacent integer endpoints. It intersects at most \(\lceil n/2\rceil+3\) of the original length-\(q\) slots. Consequently at most \(nq\,2^{\lceil n/2\rceil+3}\) of the words are reversible up to rotation. For large \(n\), choose a word \(U\) outside this set. Let \(p\) be the least period of the periodic continuation \(U^\omega\), so \(p\) divides \(|U|\). The languages of length \(p\) of \(U^\omega\) and its reversed period are disjoint, since equality of any two such words would identify the rotations and give a reflection.

Choose one of the original first returns \(D\in\{U_0,U_1\}\) with first digit different from that of \(U\), and put \(d=|D|\). Write
\[
 r=d+p+1,\qquad q_0=\left\lceil\frac{d+2p+1}{|U|}\right\rceil .
\]
Choose integers \(M_1=M_0+1\) with \(M_0|U|>4r\) and \(M_0>q_0\). From the state \(v\), follow the paths with output
\[
                  U^{M_0}D\,U^{M_{z_1}}D\,U^{M_{z_2}}D\cdots .
                                                        \tag{**}
\]
The actual digit string \(x\) includes the initial state \(v\) of length \(m\). All paths used are tight, so \eqref{eq:tight-cost} gives \(D_n=O(1)\).

The word \(V=U^{q_0}\) occurs with bounded gaps in the tail in \((**)\). Each length-\(|V|\) window there meets at most one copy of \(D\) and has more than \(2p\) digits in at most two portions of the periodic \(U\)-runs. It therefore contains \(p\) consecutive digits from \(U^\omega\). This excludes the occurrence of \(V^{\rev}\), using disjointness of the two length-\(p\) languages. A palindrome containing \(V\) would contain \(V^{\rev}\), so palindrome lengths are bounded in the tail. Symmetric trimming loses at most \(2m\) digits from a palindrome meeting the initial state.

For decoding, index the actual digits from zero and set \(f_i=1\) when \(x_i\ne x_{i+p}\) and \(f_i=0\) otherwise. Away from the initial state the nonzero flags lie within the intervals \([s-p,s+d)\) for the starts \(s\) of copies of \(D\). In each such interval the set of nonzero positions is the same translate of a fixed nonempty set \(F\). Indeed, both neighboring runs have fixed phases, and \(F=\varnothing\) would force the first digit of \(D\) to be the periodic continuation of the preceding \(U\)-run. Distinct translated sets are separated by more than \(r\) zero flags. Within one set the successive nonzero positions are at distance at most \(d+p\). Thus the marker \(c=s+\min F\) is characterized by the requirement \(f_c=1\) and \(f_{c-r}=\cdots=f_{c-1}=0\), using markers in the tail after the initial state. The first marker is at a fixed position. Its test and all later tests use only digits in \([c-r,c+p]\). Consecutive marker gaps are \(d+M_0|U|\) or \(d+M_1|U|\), and they recover \(z_1,z_2,\ldots\) in order. Put \(J_*=d+M_1|U|\).

Take a fixed integer \(B\) larger than the first marker position plus \(m+r+p\). In two equal digit intervals starting at \(t,t'\) with length \(s>B+J_*+p\), the first marker at or after \(t+B\) is at a position \(c<t+B+J_*\). Equality of the digits gives equality of the full marker tests under translation by \(t'-t\) for positions from \(c\) up to but excluding \(t+s-p\). Hence the complete gaps in this region decode identical blocks of \(z\). If the two intervals are the halves of a square of root length \(s\), the markers \(c,c+s\) are \(d_*\ge s/J_*\) gaps apart. Among the gaps from \(c\) to \(c+s\), the matching tests can lose at most \(B+J_*+p+1\) at the end, since gap lengths are at least one. We obtain a factor in \(z\) with period \(d_*\) and length at least \(2d_*-(B+J_*+p+1)\). For large \(s\) this contradicts the \(9/5\) restriction on \(z\). For two disjoint equal intervals of length \(s\), there are at least \((s-B-J_*-p)/J_*-2\) identical complete gaps on each side, disjoint also at source level. If the original intervals lie in the first \(N\) digits, the recovered gaps lie in the first \(N\) source symbols. This gives \(R_x(N)=o(N)\).

A prefix of \(N\) digits determines the first \(\lfloor (N-C_*)/J_*\rfloor\) complete gaps for a fixed \(C_*\). It reads back this many binary source symbols, and then at least \(\lfloor (N-C_*)/(21J_*)\rfloor-O(1)\) four-letter source symbols. Restrict the maximum-entropy source measure to a positive-measure subset on which all sufficiently long cylinders have mass at most \(e^{-(h-\delta)n}\), as in Lemma~\ref{lem:entropy-repetitions}. The image measure on base-\(b\) cylinders has mass \(O(e^{-(h-\delta)N/(21J_*)})\) in cylinders of length \(N\). The expansions are aperiodic since the square-root lengths are bounded. Ordinary intervals of scale \(b^{-N}\) meet at most three such cylinders. The mass distribution principle gives dimension at least \(h/(21J_*\log b)>0\) upon letting \(\delta\) decrease to zero.
\end{proof}

It follows that every path with \(\liminf D_n/n=0\) has unbounded square-root lengths exactly when \(G_{\max}\) consists of simple cycle components. The simultaneous square and palindrome conclusion holds for all such paths exactly when those components are simple cycles and all their period words are reversible up to rotation. For necessity in the simple-cycle case, the periodic path on an irreversible cycle has bounded palindrome lengths. A palindrome of length at least \(p+1\) in a period-\(p\) word would equate a full period with a rotation of its reverse. These characterizations concern all graph paths. The positive-dimensional constructions yield symbolic examples, with the algebraic-irrational case subject to further constraints.

\subsection{Explicit frequency bounds}

For binary words with palindrome lengths at most six, it suffices to forbid length-seven and length-eight palindromes. With memory seven the graph has \(112\) vertices and \(182\) edges. Table~\ref{tab:frequency-graphs} gives three rational optima. They are certified by exact comparisons of \(h(v)\ge h(u)+w(e)-\rho\) and by exhibited tight cycles. The program uses Karp's cycle-mean formula \cite{Karp} to find the candidate \(\rho\), and verifies the certificate independently of this candidate selection.

\begin{table}[ht]
\centering
\begin{tabular}{@{}llcl@{}}
\toprule
Local restriction & Event & Max. & Cyclic components\\
\midrule
Palindrome length \(\le6\) & digit \(1\) & \(2/3\) & two \(9\)-cycles\\
Palindrome length \(\le6\) & factor \(0011\) & \(1/6\) & \(18\) vertices, \(20\) edges\\
No \(000\) or \(111\) & equal neighbors & \(1/2\) & one \(4\)-cycle\\
\bottomrule
\end{tabular}
\caption{Exact frequency bounds and the components of the maximizing graph.}
\label{tab:frequency-graphs}
\end{table}

In the first row, the two period words are \(\word{110011011}\) and \(\word{110011110}\) up to rotation. Complementation gives the minimum one-frequency \(1/3\), with the same period lengths. If an algebraic irrational's binary expansion eventually has the indicated restriction, let \(A(N)\) count the ones in its original first \(N\) digits. Theorem~\ref{thm:extreme-arithmetic}, with \(P=9\), gives
\[
  A(N)-N/3,\quad 2N/3-A(N)
          \gg_\alpha \frac{(\log N)^{3/2}}{\sqrt{\log\log N}} .
\]
The cut changes each deficit by at most a fixed additive amount. For the minimum one can equivalently apply the theorem with weight zero on ones and one on zeros. If the upper density attains \(2/3\) or the lower density attains \(1/3\), Theorem~\ref{thm:extreme-frequency} gives unbounded square-root lengths. The logarithmic deficit bound still allows sublinear deficit.

In the second row the unique cyclic component has two branching vertices. At the state \(\word{0110100}\), two first returns have words \(\word{011001110100110100}\) and \(\word{110100}\). Their weights are respectively three and one, with lengths eighteen and six. Theorem~\ref{thm:branching-frequency} applies. In the last row the period \(\word{1100}\) is reversible up to rotation. Attaining upper density \(1/2\) for equal neighbors under this local restriction therefore forces both unbounded lengths in the same string.


\clearpage
\section*{Numbering correspondence}
For ease of cross-reference, the following table records the numbering of the principal statements as they appeared in an unpublished extended account of related results dated 26~November~2026.
\begin{center}
\begin{tabular}{@{}ll@{}}
\toprule
This paper & Extended account \\
\midrule
Theorem~T & Theorem~T \\
Lemma~\ref{lem:ternary} & Lemma~2.1 \\
Lemma~\ref{lem:five} & Lemma~2.2 \\
Lemma~\ref{lem:finite} & Lemma~2.3 \\
Theorem~\ref{thm:window} & Theorem~2.4 \\
Lemma~\ref{lem:pal-catalogue-cut} & Lemma~16.1 \\
Theorem~\ref{thm:six-pal} & Theorem~16.2 \\
Proposition~\ref{prop:dyadic-match} & Proposition~16.3 \\
Lemma~\ref{lem:first-gamma} & Lemma~17.1 \\
Theorem~\ref{thm:ternary-linear-match} & Theorem~17.2 \\
Theorem~\ref{thm:binary-linear-square} & Theorem~18.1 \\
Lemma~\ref{lem:entropy-repetitions} & Lemma~18.2 \\
Theorem~\ref{thm:extreme-frequency} & Theorem~19.1 \\
Theorem~\ref{thm:extreme-arithmetic} & Theorem~19.2 \\
Theorem~\ref{thm:branching-frequency} & Theorem~19.3 \\
\bottomrule
\end{tabular}
\end{center}

\begin{thebibliography}{99}
\bibitem{ABL}
B. Adamczewski, Y. Bugeaud, and F. Luca,
\emph{Sur la complexit\'e des nombres alg\'ebriques},
C. R. Math. Acad. Sci. Paris 339 (2004), 11--14.
\href{https://doi.org/10.1016/j.crma.2004.04.012}{DOI}.

\bibitem{AR}
B. Adamczewski and N. Rampersad,
\emph{On patterns occurring in binary algebraic numbers},
Proc. Amer. Math. Soc. 136 (2008), 3105--3109.
\href{https://doi.org/10.1090/S0002-9939-08-09319-2}{DOI}.
\href{https://adamczewski.perso.math.cnrs.fr/Overlap.pdf}{Author version}.

\bibitem{BC}
G. Badkobeh and M. Crochemore,
\emph{Fewest repetitions in infinite binary words},
RAIRO Theor. Inform. Appl. 46 (2012), 17--31.
\href{https://doi.org/10.1051/ita/2011109}{DOI}.

\bibitem{BO}
G. Badkobeh and P. Ochem,
\emph{Characterization of some binary words with few squares},
Theoret. Comput. Sci. 588 (2015), 73--80.
\href{https://doi.org/10.1016/j.tcs.2015.03.044}{DOI}.
\href{https://www.lirmm.fr/~ochem/morphisms/avoidability.pdf}{Author version}.

\bibitem{Bugeaud}
Y. Bugeaud,
\emph{Distribution Modulo One and Diophantine Approximation},
Cambridge Tracts in Mathematics, vol.~193, Cambridge University Press,
2012.
\href{https://doi.org/10.1017/CBO9781139017732}{DOI}.

\bibitem{OR}
P. Ochem and M. Rosenfeld,
\emph{Avoidability of palindrome patterns},
Electron. J. Combin. 28 (2021), no.~1, Paper P1.4.
\href{https://doi.org/10.37236/9593}{DOI}.

\bibitem{SS}
J. Shallit and A. Shur,
\emph{Subword complexity and power avoidance},
Theoret. Comput. Sci. 792 (2019), 96--116.
\href{https://doi.org/10.1016/j.tcs.2018.09.010}{DOI}.

\bibitem{Bilu}
Y. F. Bilu, \emph{The many faces of the subspace theorem [after
Adamczewski, Bugeaud, Corvaja, Zannier...]}, S\'eminaire Bourbaki,
expos\'e 967, Ast\'erisque 317 (2008), 1--38.
\href{https://www.numdam.org/book-part/AST_2008__317__1_0/}{Publication record}.
\href{https://arxiv.org/pdf/0907.2098#page=8}{Theorem 3.3, printed p.~8}.

\bibitem{BE}
Y. Bugeaud and J.-H. Evertse, \emph{On two notions of complexity of algebraic numbers}, Acta Arith. 133 (2008), 221--250. \href{https://arxiv.org/abs/0709.1560}{arXiv:0709.1560}.

\bibitem{EF}
J.-H. Evertse and R. G. Ferretti, \emph{A further improvement of the Quantitative Subspace Theorem}, Ann. of Math. 177 (2013), 513--590. \href{https://doi.org/10.4007/annals.2013.177.2.4}{DOI}.

\bibitem{Hall}
The Hall--Thue fixed point and its square-free property, as stated in \href{https://doi.org/10.1016/j.tcs.2022.04.025}{Avoiding square-free words on free groups} (2022). See also the Hall-word discussion in \href{https://www.lirmm.fr/~ochem/morphisms/pairs.pdf}{Theoret. Comput. Sci. 799 (2019), 140--148}.

\bibitem{KS}
J. Karhum\"aki and J. Shallit, \emph{Polynomial versus exponential growth in repetition-free binary words}. \href{https://arxiv.org/abs/math/0304095}{arXiv:math/0304095}.

\bibitem{RSW}
N. Rampersad, J. Shallit, and M.-w. Wang, \emph{Avoiding large squares in infinite binary words}. \href{https://arxiv.org/abs/math/0306081}{arXiv:math/0306081}.

\bibitem{Stephan}
R. Stephan,
\emph{The Evertse--Ferretti interval result for a set of weights, and the exponent $1/2$ in three complexity bounds},
preprint (2026).
\href{https://arxiv.org/abs/2611.16536v1}{arXiv:2611.16536v1}.

\bibitem{Breiman}
L. Breiman,
\emph{The individual ergodic theorem of information theory},
Ann. Math. Statist. 28 (1957), 809--811.
\href{https://doi.org/10.1214/aoms/1177706899}{DOI}.

\bibitem{CurrieTest}
J. D. Currie,
\emph{Finite test sets for morphisms that are squarefree on some of Thue's squarefree ternary words},
J. Integer Seq. 22 (2019), Article 19.8.2.
\href{https://cs.uwaterloo.ca/journals/JIS/VOL22/Currie/currie20.html}{Journal text}.

\bibitem{Dejean}
F. Dejean, \emph{Sur un th\'eor\`eme de Thue},
J. Combin. Theory Ser. A 13 (1972), 90--99.
\href{https://doi.org/10.1016/0097-3165(72)90011-8}{DOI}.

\bibitem{DMO}
L'. Dvo\v r\'akov\'a, L. Mol, and P. Ochem,
\emph{Critical exponent of ternary words with few distinct palindromes},
Discrete Math. Theor. Comput. Sci. (2026).
\href{https://doi.org/10.46298/dmtcs.16125}{DOI}.
\href{https://arxiv.org/abs/2507.04925v3}{Version 3}.
\href{https://www.lirmm.fr/~ochem/morphisms/few_ternary_palindromes.pdf}{Authors' copy}.
\href{https://www.lirmm.fr/~ochem/morphisms/palin3.htm}{Authors' supplement}.

\bibitem{DOO}
L'. Dvo\v r\'akov\'a, P. Ochem, and D. Opo\v censk\'a,
\emph{Critical exponent of binary words with few distinct palindromes},
Electron. J. Combin. 31 (2024), Paper P2.29.
\href{https://doi.org/10.37236/12574}{DOI}.
\href{https://arxiv.org/abs/2311.13003v2}{Version 2}.

\bibitem{Karp}
R. M. Karp, \emph{A characterization of the minimum cycle mean in a digraph},
Discrete Math. 23 (1978), 309--311.
\href{https://doi.org/10.1016/0012-365X(78)90011-0}{DOI}.

\bibitem{OchemSource}
P. Ochem, \emph{A generator of morphisms for infinite words},
RAIRO Theor. Inform. Appl. 40 (2006), 427--441.
\href{https://doi.org/10.1051/ita:2006020}{DOI}.


\end{thebibliography}

\end{document}
